\documentclass[11pt]{article}
\usepackage[margin=1in]{geometry}
\usepackage{amsmath,amssymb}
\usepackage{booktabs}
\usepackage{graphicx}
\usepackage{xcolor}
\usepackage[hidelinks]{hyperref}
\usepackage{natbib}
\usepackage{caption}
\usepackage{algorithm}
\usepackage{algpseudocode}
\usepackage{amsthm}
\newtheorem{proposition}{Proposition}
\newtheorem{corollary}{Corollary}

\newcommand{\jev}{\textsc{Jev}}
\newcommand{\todo}[1]{\textcolor{red}{[#1]}}

\title{Typed, Calibrated Decisions as a Reward Source \emph{and} as a Critic\\
for Reinforcement Learning Fine-Tuning of Language Models}

\author{%
  \\ \\
  \texttt{}%
}
\date{}

\begin{document}
\maketitle

\begin{abstract}
Reinforcement learning fine-tuning of language models needs a source of reward,
and every source in common use is a compromise. A Bradley--Terry reward model is
cheap but frozen in the preference distribution it was fit to and capped by a
short context window. An LLM judge is flexible but expensive, slow, and read
through a lossy text-parsing step. A self-judge is free but, at small scale,
close to noise. We study a fourth option that has not previously been used for
RL: a \emph{System One} model, which emits typed, calibrated answers to
predeclared questions and never generates text. We use \jev{} (TypeSafe) in two
roles. As a \textbf{reward function} it answers a fixed rubric about a finished
response in one request, returning per-question probabilities directly, with no
parsing. As a \textbf{critic} it answers a forward-looking question about a
\emph{partial} response --- ``will this, once finished, be correct and
helpful?'' --- which is exactly the quantity PPO's value network is trained to
predict. We show empirically that this prefix probability tracks eventual
response quality
(AUC $=0.86$ at the full response and $=0.76$ after only a quarter of it; $R^2=0.56$
\unskip),
which licenses replacing PPO's learned value head with a frozen, zero-shot,
externally calibrated critic: no value parameters, no value loss, no
bootstrapping from a critic that is itself still learning.
We compare five reward sources under GRPO at a matched reward-call budget, and
two critics under PPO, holding the start checkpoint, prompt set, prompt order,
decoding parameters and KL coefficient fixed, and evaluating at matched
KL rather than at matched step count. We report a cross-reward matrix that
separates genuine improvement from over-optimisation of whichever judge a run
was trained against, and full cost accounting.
\todo{headline results once runs complete}
\end{abstract}

\section{Introduction}

The standard recipe for aligning a language model with reinforcement learning
fixes an optimiser --- PPO \citep{schulman2017ppo,ouyang2022instructgpt} or a
critic-free variant such as GRPO \citep{shao2024deepseekmath} --- and varies
almost everything else. In practice the component that determines whether the
run is useful, and most of what it costs, is the \emph{reward source}: the thing
that turns a sampled response into a number.

Three families dominate. A \textbf{learned reward model} fits a Bradley--Terry
objective to human preference pairs \citep{bradley1952rank,stiennon2020summarize}
and is then frozen; it is nearly free to query, but inherits the distribution of
the preference data and, for the encoder models typically used, a 512-token
window. \textbf{LLM-as-judge} / RLAIF prompts a capable instruct model to grade
the response \citep{bai2022constitutional,lee2024rlaif,zheng2023judging}; it
generalises far better, but each judgment is an autoregressive generation that
must be parsed, and the parse is a lossy, brittle step between the judge's
belief and the scalar the optimiser sees. \textbf{Self-rewarding} setups
\citep{yuan2024selfrewarding} let the policy grade itself, which removes the
external cost but, at the model scales most practitioners can afford to train,
produces a judge barely better than chance.

A fourth kind of model has recently become available and does not fit this
taxonomy. TypeSafe's \jev{} \citep{typesafe2026jev} is a \emph{System One} model:
it is given a state and a set of questions declared in advance, and it returns
\emph{typed} answers --- a probability for a yes/no proposition, a distribution
over named options, an expected level on an ordered scale --- together with
calibration information. It does not produce text. Because the answer space is
fixed by the request, the model cannot return a value outside the schema, and
because no tokens are decoded, a judgment costs roughly one prefill and returns
in a fraction of a second.

For reinforcement learning this combination is unusually well matched to the
problem, in two separate ways, and the distinction between them is the subject
of this paper.

\paragraph{As a reward function.} A rubric is a set of predeclared questions.
That is precisely \jev{}'s input format, so a multi-criterion reward becomes one
request, the per-criterion answers come back as numbers rather than as prose to
be parsed, and the scalarisation is the experimenter's choice rather than an
artefact of how a judge happened to phrase itself.

\paragraph{As a critic.} In a language-model MDP with a single terminal reward
$R$, the value function is
\begin{equation}
V^{\pi}(s_t) \;=\; \mathbb{E}\!\left[\, R \,\middle|\, x,\, y_{<t} \,\right],
\label{eq:value}
\end{equation}
the expected quality of the finished response given the prefix generated so far.
This is a question about a partial state with a yes/no answer, which is exactly
the shape \jev{} accepts. PPO normally learns \eqref{eq:value} with a value head
trained online against its own returns. If an external model can answer it
zero-shot, in calibrated units, for a fraction of a cent, the value network
becomes unnecessary --- and with it the value parameters, the value loss, the
value optimiser state, and the bootstrapping bias of a critic that is itself
still being fit. To our knowledge no prior work uses an external judge as a
token-level critic; judges are used to produce terminal rewards.

\paragraph{A design note on what we did \emph{not} test.} A natural fourth
condition --- ``\jev{} as an evaluator'' --- collapses into the reward-function
condition. Scoring a finished response and taking the resulting scalar as the
learning signal \emph{is} using \jev{} as a reward function; running it as a
separate arm would duplicate a column. We therefore drop it and spend the
comparison budget on the critic role, which is genuinely distinct because it
queries states the reward function never sees.

\paragraph{Contributions.}
\begin{enumerate}
  \item The first study of a typed-decision (System One) model as a reward
    source for RL fine-tuning, compared against a Bradley--Terry reward model,
    an LLM judge (RLAIF) and a self-judge under an identical rubric, identical
    scalarisation, and a matched budget of reward calls.
  \item \jev{} as a \emph{critic}: we show its prefix probability is an
    empirically calibrated estimator of final response quality, and use it to
    run PPO with no value network at all.
  \item A fairness protocol that reads the LLM and self-judges
    \emph{from their logits} rather than by parsing text, so that the
    dense, parse-free property is not silently granted to \jev{} alone.
  \item An ablation on where the rubric comes from
    (\textsc{StructuredGeneration2Jev}): fixed, generated per prompt by a frozen
    LLM, or proposed by the policy being trained --- the last as a deliberate
    reward-hacking probe.
  \item Cost and latency accounting for every arm, measured rather than quoted.
\end{enumerate}

\section{Background}

\subsection{RL fine-tuning}
We use two optimisers. \textbf{PPO} \citep{schulman2017ppo} maximises a clipped
surrogate with advantages from GAE \citep{schulman2016gae}, requiring a critic
$V$. \textbf{GRPO} \citep{shao2024deepseekmath} removes the critic by sampling
$G$ completions per prompt and using the group mean as the baseline; it is the
natural setting in which to vary the reward source alone, since no value network
confounds the comparison.

\subsection{Reward sources}
\label{sec:sources}
See Section~\ref{sec:method} for exact definitions. Briefly: a frozen DeBERTa
Bradley--Terry model; a frozen instruct LLM read through logits; a frozen copy of
the SFT checkpoint read the same way; \jev{}; and, as a degenerate control, the
policy's own token-level self-certainty, which consults nothing external.

\subsection{System One models and \jev{}}
\jev{} accepts a \emph{state} and a map of \emph{questions}, and returns one
typed answer each \citep{typesafe2026jev}. Three question types exist: a
\textsc{noul} returns $P(\text{true}) \in [0,1]$ for a proposition; a
\textsc{choice} returns a distribution over named options plus a confidence; a
\textsc{score} returns an expected level on an ordered rubric of $2$--$10$
levels, plus the level distribution and a confidence. All questions in a request
are answered against the same state, and adding questions has little effect on
latency. We access \jev{} 1.13 through OpenRouter's decisions endpoint. Measured
properties relevant to RL are given in Table~\ref{tab:jevprops}.

\begin{table}[t]\centering
\caption{Measured properties of \jev{} 1.13.}\label{tab:jevprops}
\textcolor{red}{[pending: run not yet complete]}
\end{table}


\section{Method}
\label{sec:method}

\subsection{A single rubric, rendered into each judge's native format}
\label{sec:rubric}
Every judge-based arm answers the same five questions about a
(request, response) pair: \textsc{correct} (noul, $+$), \textsc{helpful}
(3-level score, $+$), \textsc{padded} (noul, $-$), \textsc{evasive} (noul, $-$)
and \textsc{self\_promo} (noul, $-$). Writing $\hat{q}_i \in [0,1]$ for the
normalised answer to question $i$ (a score of $L$ levels is divided by $L-1$),
the reward is
\begin{equation}
R \;=\; \hat{q}_{\text{correct}} \;+\; \hat{q}_{\text{helpful}}
      \;-\; \hat{q}_{\text{padded}} \;-\; \hat{q}_{\text{evasive}}
      \;-\; \hat{q}_{\text{self\_promo}} \;\in\; [-3, +2].
\label{eq:reward}
\end{equation}
\textsc{self\_promo} is not optional. The policy's own text sits inside the
judge's state, so a policy can learn to assert its own quality at the judge; we
both penalise it in the reward and track it as a hacking detector
(Section~\ref{sec:eval}).

\subsection{Reading judges without parsing}
\label{sec:parsefree}
Comparing a typed model against a text-generating judge invites an unfair
advantage: if the LLM judge is read by parsing its prose, we are partly
measuring parse failures. We therefore read every generative judge from its
logits. For a noul we take
$P(\textsf{yes}) / (P(\textsf{yes}) + P(\textsf{no}))$, summing over
capitalisation and leading-space variants of each token; for a score we take a
softmax over the level-index tokens and its expectation. Each question is a
separate prefill sharing the state prefix, mirroring how \jev{} evaluates
questions in parallel. The scalarisation \eqref{eq:reward} is then byte-identical
across \jev{}, the RLAIF judge and the self-judge.

This applies to the \emph{hosted} RLAIF judge as well, which matters because
the judges small enough to run on our hardware sit at or near chance on held-out
preference pairs (Section~\ref{sec:res-judges}); comparing against one of those
would be comparing against a straw man. Three practical details make a hosted
judge readable:

\begin{itemize}
  \item \textbf{Reasoning must be off.} A reasoning model spends its first
    tokens on thought and returns no content to read; with reasoning disabled
    the answer token is the first token.
  \item \textbf{The provider must be pinned.} The gateway load-balances across
    providers and only some return \texttt{logprobs}; without requiring the
    parameter, a fraction of calls silently fall back to parsing. Requiring it
    took the unreadable fraction from 50\% to under 5\%.
  \item \textbf{Prompt order is state-then-question.} The rubric's questions
    about one completion share a long prefix (system and state) and differ only
    in a short suffix, so the state goes first. This is the reverse of the usual
    ``instructions first'' caching advice, which optimises the wrong axis when
    the repeated span is the state rather than the instruction.
\end{itemize}

Calls whose logprobs are still unreadable fall back to the emitted word, and we
report that fraction rather than leaving it implicit.

\subsection{\jev{} as a reward function}
One request carries the state and all five rubric questions; the answers are
normalised and combined by \eqref{eq:reward}. No parsing, no retry-on-malformed,
and one network round trip per completion regardless of rubric size.

\subsection{\jev{} as a critic}
\label{sec:critic}
We ask \jev{} a single forward-looking noul about a \emph{partial} response:
\emph{if the assistant continues in the same direction and finishes, will the
completed response be a correct and genuinely helpful answer?} Write $p_t$ for
the returned probability at prefix length $t$.

\paragraph{Calibration.} $p_t$ lives in $[0,1]$ and $R$ in $[-3,+2]$, so TD
errors computed from raw $p_t$ would be in arbitrary units. We fit
$\hat{V}(s_t) = a\,p_t + b$ by least squares on held-out
(complete response, reward) pairs drawn from the SFT policy, and report $R^2$,
so the quality of the map is visible rather than assumed
(Section~\ref{sec:res-calibration}).

\paragraph{Sparse probing.} Querying every token position is unnecessary and
would dominate cost. We probe $K$ positions per completion, evenly spaced in
token index, and interpolate $\hat{V}$ piecewise-linearly between them; $K$ is an
ablation axis. Advantages then follow standard GAE with $\gamma = 1$, using
$\hat{V}(s_{T}) = 0$ at the terminal state.

\paragraph{What this removes.} The Jev-critic arm has no value head, no value
loss term, no value learning rate and no value optimiser state. Against the
learned-critic arm it differs in exactly that, since both run through the same
PPO implementation.

\begin{algorithm}[t]
\caption{PPO with \jev{} as the critic. Differences from textbook PPO are the
lines marked $\triangleright$; everything else, including the clipped surrogate
and the KL shaping, is unchanged.}
\label{alg:jevcritic}
\begin{algorithmic}[1]
\State \textbf{calibrate once:} sample $(x, y)$ from $\pi_{\text{SFT}}$;
  fit $a, b$ minimising $\sum (a\,p(x,y) + b - R(x,y))^2$
  \Comment{$\triangleright$ puts $\hat V$ in reward units}
\For{each step}
  \State sample completions $y \sim \pi_\theta(\cdot \mid x)$
  \State $R \gets$ rubric$(x, y)$ \Comment{terminal reward, one \jev{} call per completion}
  \State $r_t \gets -\beta\big(\log \pi_\theta(y_t) - \log \pi_{\text{ref}}(y_t)\big)$;\;
         $r_{T-1} \mathrel{+}= R$
  \For{$k = 1 \dots K$} \Comment{$\triangleright$ $K$ prefix probes, not a value net}
    \State $t_k \gets \lfloor k(T-1)/(K-1) \rceil$;\quad
           $p_{t_k} \gets \textsc{Jev}\big(x,\, y_{<t_k},\, q_{\text{prefix}}\big)$
  \EndFor
  \State $\hat V(s_t) \gets a \cdot \textsc{interp}(t; \{t_k\}, \{p_{t_k}\}) + b$,\;
         $\hat V(s_T) \gets 0$
         \Comment{$\triangleright$ frozen, no gradient}
  \State $\delta_t \gets r_t + \gamma \hat V(s_{t+1}) - \hat V(s_t)$;\;
         $A_t \gets \sum_{l} (\gamma\lambda)^{l} \delta_{t+l}$
  \State update $\theta$ on the clipped surrogate
         \Comment{$\triangleright$ no value loss term}
\EndFor
\end{algorithmic}
\end{algorithm}

\subsection{What sparse probing costs in principle}
\label{sec:analysis}

Probing $K$ positions and interpolating is an approximation, and it is worth
saying exactly what it approximates and how badly, because that is what decides
how to choose $K$ and what the resulting advantages mean.

\begin{proposition}[Credit is assigned per segment, not per token]
\label{prop:segments}
Let $\hat V$ be affine on each interval between consecutive probes
$t_k, t_{k+1}$, and let $\gamma = 1$. Then for $t_k \le t < t_{k+1} - 1$ the
temporal-difference error is
\[
\delta_t \;=\; r_t + \hat V(s_{t+1}) - \hat V(s_t)
\;=\; r_t + \frac{p_{t_{k+1}} - p_{t_k}}{t_{k+1} - t_k}\,a,
\]
which is constant in $t$ within the segment apart from the per-token shaping
term $r_t$. Consequently the critic distinguishes credit only at the
granularity of the $K-1$ segments.
\end{proposition}

This is a structural difference from a learned value head rather than an
artefact of insufficient probing: no increase in per-probe accuracy makes the
interpolated critic discriminate between two tokens inside one segment. It is
also why $K$ is an experimental axis rather than an implementation detail.

\begin{proposition}[Bias of the interpolated critic]
\label{prop:bias}
Write $V^\ast(t)$ for the true value at token index $t$ and $h = T/(K-1)$ for
the probe spacing. If $V^\ast$ is $L$-Lipschitz in $t$, then
\[
\sup_t \big| \hat V_K(t) - V^\ast(t) \big| \;\le\; \tfrac{1}{2} L h
\;=\; O\!\big(T/K\big).
\]
If in addition the second difference of $V^\ast$ is bounded by $M$, the bound
tightens to $\tfrac{1}{8} M h^2 = O\!\big(T^2/K^2\big).$
\end{proposition}

Both are the standard linear-interpolation bounds; what is not standard is
which regime a real value curve occupies, and that is an empirical question we
answer in Section~\ref{sec:res-density} by probing every position on a small
sample and measuring how the error actually decays with $K$.

\begin{corollary}[The bias reaches the optimiser undiminished]
\label{cor:adv}
With $\gamma = \lambda = 1$ the generalised advantage telescopes to
$A_t = R - \hat V(s_t)$, so the advantage error equals the value error
pointwise and inherits the bounds of Proposition~\ref{prop:bias}.
\end{corollary}

\paragraph{Cost of the critic.} The arm spends $K$ extra \jev{} calls per
completion. At the measured price this is a fraction of a cent per step and
buys token-level credit assignment; the comparison in
Section~\ref{sec:res-calibration} is what decides whether the signal is worth
it, and Table~\ref{tab:cost} reports what it cost.

\subsection{\textsc{StructuredGeneration2Jev}}
\label{sec:sg2jev}
The component that decides \emph{which} questions \jev{} is asked is a swappable
axis with three settings. \textbf{static}: the fixed rubric of
Section~\ref{sec:rubric}; the only setting under which \jev{}, the LLM judge and
the self-judge are comparable, and therefore the setting used in the main
pipeline. \textbf{generated}: a frozen LLM writes up to three task-specific
questions per \emph{prompt}, each carrying an explicit sign, appended to the
fixed rubric; questions are generated per prompt and cached, never per
completion, since a reward that depends on the completion through the question
set is not a function of the completion and makes the RL objective ill-posed.
\textbf{policy}: the model being trained proposes its own questions --- a
deliberate reward-hacking probe, since nothing prevents it from proposing
criteria it trivially satisfies. Generated questions are schema-validated and
dropped if malformed; rewards are rescaled by the rubric's attainable range so
that a longer rubric does not inflate the reward scale.

\subsection{Controlled variables}
\label{sec:controls}
Identical across every run: the SFT start checkpoint, the prompt set and its
order (fixed dataloader seed, verified by a logged fingerprint), decoding
parameters, KL coefficient, $G$ for GRPO, and the \emph{budget of reward calls}
rather than the number of steps. We use three seeds per arm and report
mean $\pm$ standard deviation.

\section{Experimental setup}

The policy is a small Qwen2.5 instruct model \citep{qwen2025qwen25}; prompts and
preference pairs come from UltraFeedback \citep{cui2023ultrafeedback}; training
uses TRL \citep{vonwerra2020trl} for SFT, GRPO and the reward model, and our own
PPO implementation for the critic comparison, because TRL 1.13 no longer ships a
PPO trainer and because one implementation with a swappable critic is what makes
that comparison exact.

\begin{table}[t]\centering\small
\caption{Configuration. Every entry is held fixed across all arms except the
reward source (GRPO phase) and the critic (PPO phase).}
\label{tab:setup}
\begin{tabular}{ll}
\toprule
Component & Setting \\
\midrule
Policy & \texttt{Qwen/Qwen2.5-0.5B-Instruct} \\
SFT data & UltraFeedback (binarized), chosen responses \\
RL prompts & UltraFeedback \texttt{train\_prefs}, fixed order \\
Held-out prompts & UltraFeedback \texttt{test\_prefs}, 200 \\
Ground-truth task & GSM8K, 250 test problems \\
BERT RM & \texttt{OpenAssistant/reward-model-deberta-v3-large-v2} \\
RLAIF judge & \texttt{Qwen/Qwen2.5-1.5B-Instruct}, read from logits \\
Self-judge & frozen copy of the SFT checkpoint \\
\jev{} & \texttt{typesafe/jev-1.13} \\
Evaluation judge & \texttt{openai/gpt-4.1-mini} (unrelated to every training judge) \\
GRPO & $G=8$, 16 prompts/step, 250 steps, lr $10^{-6}$, $\beta=0.04$ \\
PPO & 64 prompts/step, 150 steps, lr $10^{-6}$, KL coef 0.05, GAE $\lambda=0.95$ \\
Reward-call budget & 32,000 per GRPO run (matched across sources) \\
Decoding & temperature 1.0, top-$p$ 1.0, max 384 new tokens \\
Seeds & 3 per core arm \\
Hardware & 2$\times$NVIDIA T4 (16\,GB), Kaggle \\
\bottomrule
\end{tabular}
\end{table}


\paragraph{Why the policy is adapted rather than fully fine-tuned.} On the
14.6\,GB card available, full fine-tuning of even a 0.5B policy does not fit at
\emph{any} micro-batch size, and the reason is worth stating because it is not
the one usually assumed. The weights are 1.8\,GB; the gradients are another
1.8\,GB, the optimiser moments 3.7\,GB, the separate reference policy 1.8\,GB
and the rollout KV cache 2.6\,GB --- together over the card before a single
logit is allocated, and the logits for a 150k vocabulary are themselves
4\,GB at a micro-batch of eight. LoRA removes three of those terms at once,
since the reference policy becomes the base model with adapters disabled.
Every arm uses an identical adapter configuration, so the comparison between
reward sources is unaffected; absolute quality under full fine-tuning would
differ, and we do not claim otherwise. Our memory model
(\texttt{scripts/memory\_budget.py}) reproduces the observed failure to within
0.5\,\% and is included so the calculation can be redone for other hardware.

\subsection{Evaluation}
\label{sec:eval}
\paragraph{Quality.} Head-to-head win-rate against the SFT checkpoint on held-out
prompts, judged by a model from a family used by \emph{no} training reward
source. Each pair is judged twice with positions swapped and averaged, removing
position bias. The headline number is length-controlled in the AlpacaEval-2
sense \citep{dubois2024lengthcontrolled}: we fit
$P(\text{win}) = \sigma(\theta + \gamma\,\Delta\ell)$ over standardised length
differences and report $\sigma(\theta)$. Confidence intervals are bootstrap over
prompts.

\paragraph{Matched divergence.} Arms are compared at equal
$\mathrm{KL}(\pi \,\|\, \pi_{\text{ref}})$, not at equal step count; a reward
source that moves faster per step is not thereby better.

\paragraph{Hacking detectors.} GSM8K accuracy \citep{cobbe2021gsm8k} (a judge can be talked into
approval; arithmetic cannot), mean response length, self-praise phrase rate,
distinct-2, and the gap between training reward and held-out quality.

\paragraph{Cross-reward matrix.} Over-optimising a learned proxy is a
well-documented failure of RLHF \citep{gao2023scaling}, and a judge-based reward
is a proxy like any other. Every final checkpoint is scored by
\emph{every} reward source. The diagonal is what a run optimised; the
off-diagonal is what the other judges make of it. We report
$\text{diag} - \overline{\text{off-diag}}$ in baseline-relative units as an
over-optimisation index.

\paragraph{Cost.} Dollars, wall-clock, reward-call count and reward latency
(p50/p95), measured in-run rather than quoted from documentation.

\section{Results}

\subsection{Can each reward source recognise a human preference at all?}
\label{sec:res-judges}
A comparison is only informative if every arm's judge works. Table~\ref{tab:judges}
measures each source on held-out preference pairs before any RL, which is what
fixes the RLAIF judge size; running that arm with a judge at chance would be a
straw man rather than a baseline.
\begin{table}[t]\centering\small
\caption{How well each reward source recognises a human preference, measured
before any RL. Each source scores the chosen and the rejected response of
$n=1400$ held-out UltraFeedback pairs; accuracy is how often it ranks the
chosen one higher, and Cohen's $d$ is the separation in units of the margin's
own standard deviation. A source that cannot do this will not teach a policy
anything, so this is what fixes the RLAIF judge size rather than an assumption.
Wall-clock is for all $2n$ scorings on the same hardware.}
\label{tab:judges}
\begin{tabular}{lcccc}
\toprule
Reward source & Accuracy (\%) & Cohen's $d$ & Mean margin & Wall-clock (s) \\
\midrule
api\_judge:deepseek-v4-flash-0731 & 68.0 & +0.53 & +0.573 & 539 \\
\jev{} & 66.3 & +0.43 & +0.487 & 19 \\
LLM judge Qwen2.5-7B-Instruct & 65.1 & +0.40 & +0.506 & 8704 \\
LLM judge Qwen2.5-3B-Instruct & 62.3 & +0.32 & +0.239 & 3767 \\
BERT RM (DeBERTa BT) & 58.8 & +0.22 & +0.399 & 162 \\
\bottomrule
\end{tabular}
\end{table}


Two things in that table matter more than the ranking. First, the local judges
that fit on the available hardware are weak: the 1.5B model is at chance, and
the 3B model is below the Bradley--Terry reward model it is supposed to
improve on. That is why the RLAIF arm uses a hosted judge, read from logits so
that the fairness protocol of Section~\ref{sec:parsefree} still holds.

Second, and this is the efficiency claim the paper rests on:
\jev{} scores 300 pairs at 0.640 accuracy in 6 seconds, while the 7B local
judge reaches 0.647 in 1831 seconds --- a difference of 0.7 accuracy points
against a factor of 310 in wall-clock. The 7B judge also had to be sharded
across both GPUs to run at all, occupying roughly 22\,GB, against none for the
typed model. We do not claim \jev{} is the better judge; at this sample size
those two accuracies are not distinguishable, which is precisely the point.
The claim is that equivalent judgement is available at a cost that makes
querying it $K$ times per completion --- the critic role --- affordable, and
that this is what changes which algorithms are reachable.

\subsection{Is \jev{}'s prefix probability a value function?}
\label{sec:res-calibration}
\begin{table}[t]\centering\small
\caption{\jev{}'s forward-looking probability on a \emph{partial} response,
as a function of how much of the response it has seen, measured on
$n=256$ completions sampled from the SFT policy. ``Good'' and ``bad'' are
the above- and below-median halves by \emph{final} rubric reward, which is not
observable at the time of the query. AUC is the probability that the prefix
value ranks an eventually-good response above an eventually-bad one; $R^2$ is
for the affine fit of final reward on prefix value.}
\label{tab:calib}
\begin{tabular}{lccccc}
\toprule
Prefix seen & $\bar{p}$ (good) & $\bar{p}$ (bad) & AUC & Spearman $\rho$ & $R^2$ \\
\midrule
0\% & 0.539 & 0.506 & 0.608 & 0.222 & 0.029 \\
10\% & 0.619 & 0.476 & 0.658 & 0.318 & 0.105 \\
25\% & 0.574 & 0.347 & 0.754 & 0.487 & 0.207 \\
50\% & 0.498 & 0.239 & 0.784 & 0.556 & 0.303 \\
75\% & 0.453 & 0.184 & 0.793 & 0.593 & 0.374 \\
100\% & 0.452 & 0.143 & 0.802 & 0.635 & 0.520 \\
\bottomrule
\end{tabular}
\end{table}


\subsection{How fast does the sparse-probe bias vanish?}
\label{sec:res-density}
\begin{table}[t]\centering
\caption{Interpolation error against the number of probes.}\label{tab:density}
\textcolor{red}{[pending: run not yet complete]}
\end{table}


The mean error decays as $K^{-1.03}$, which is the Lipschitz rate of
Proposition~\ref{prop:bias} essentially exactly, and \emph{not} the
$K^{-2}$ that bounded curvature would give. The worst-case error behaves
differently: it decays as $K^{-0.35}$, far more slowly, and is still
0.12 at $K=16$.

We read the gap between the two as evidence that the value curve is Lipschitz
but not smooth --- that it contains localised near-discontinuities. That is what
one should expect from the task rather than from the estimator: a response can
commit to being bad at a particular token, by opening a self-promotional
clause or by turning evasive, and the expected final quality drops abruptly
there. No amount of uniform probing recovers a step.

Two things follow. First, $K$ buys mean accuracy efficiently and worst-case
accuracy barely, so the usual instinct to raise $K$ until the critic "looks
converged" spends calls on the wrong quantity. At $K=6$, the setting we use,
the mean error is 0.036 in probability units, which the calibration of
Section~\ref{sec:res-calibration} maps to roughly 0.08 in reward units on a
scale of width 5. Second, and more usefully, it suggests that probes should not
be placed uniformly: allocating them where the value curve moves fastest would
attack precisely the error that $K$ does not. We leave adaptive probe placement
to future work, but note that the measurement here is what identifies it as the
right thing to try.

\subsection{GRPO: varying the reward source}
\begin{table}[t]\centering\small
\caption{GRPO with the reward source as the only manipulated variable, at a matched budget of reward calls. Win-rate is against the SFT checkpoint, length-controlled, judged by a model unrelated to every training judge. Mean $\pm$ s.d.\ over three seeds.}
\label{tab:grpo}
\begin{tabular}{lcccccc}
\toprule
Reward source & LC win-rate (\%) & Raw (\%) & GSM8K (\%) & Len (tok) & Self-praise (\%) & KL/tok \\
\midrule

\bottomrule
\end{tabular}
\end{table}


\subsection{PPO: varying the critic}
\begin{table}[t]\centering\small
\caption{PPO with the critic as the only manipulated variable: a learned value head versus \jev{} valuing token prefixes zero-shot. Both arms use the same reward source and the same PPO implementation.}
\label{tab:ppo}
\begin{tabular}{lcccccc}
\toprule
Critic & LC win-rate (\%) & Raw (\%) & GSM8K (\%) & Len (tok) & Self-praise (\%) & KL/tok \\
\midrule

\bottomrule
\end{tabular}
\end{table}


\subsection{Reward hacking and over-optimisation}
\begin{table}[t]\centering
\caption{Cross-reward matrix.}\label{tab:cross}
\textcolor{red}{[pending: run not yet complete]}
\end{table}


\subsection{Where the rubric comes from}
\begin{table}[t]\centering\small
\caption{Where the \jev{} question set comes from (\textsc{StructuredGeneration2Jev}). All three arms are otherwise identical GRPO runs with \jev{} as the reward source. The \emph{policy} arm lets the model being trained write its own grading criteria and is included as a reward-hacking probe.}
\label{tab:sg2jev}
\begin{tabular}{lcccccc}
\toprule
Rubric source & LC win-rate (\%) & Raw (\%) & GSM8K (\%) & Len (tok) & Self-praise (\%) & KL/tok \\
\midrule

\bottomrule
\end{tabular}
\end{table}


\subsection{Cost}
\begin{table}[t]\centering\small
\caption{Measured cost of one training run per reward source. Local judges cost
GPU-hours rather than dollars; API judges cost both. Latency is per reward call
as observed by the trainer, including network round-trip.}
\label{tab:cost}
\begin{tabular}{lccccc}
\toprule
Arm & Reward calls & USD & p50 (s) & p95 (s) & GPU-h \\
\midrule

\bottomrule
\end{tabular}
\end{table}


\section{Discussion}

\paragraph{What a typed interface buys, and what it does not.} The advantages
we can attribute to \jev{} are structural rather than mysterious: a rubric is a
set of predeclared questions, which is its native input; answers arrive as
numbers, so nothing is lost to parsing; and a judgment costs one prefill, so a
dense multi-criterion reward is affordable at RL scale. None of this implies
\jev{} is a \emph{better} evaluator than a large LLM judge, and
Section~\ref{sec:res-judges} measures that separately rather than assuming it.
The interesting claim is narrower: for the specific shape of question RL needs
answered, millions of times, the typed interface removes cost and failure modes
that the generative interface imposes.

\paragraph{Why the critic role is the one that could not be filled before.}
An LLM judge could in principle be asked our prefix question too. What makes it
impractical is arithmetic: a critic is queried at $K$ positions per completion,
so the query count is $K$ times the reward count, and at LLM-judge latency and
price that dominates the run. The typed model changes the constant factor by
enough to change which algorithms are reachable. This is the sense in which the
result is about a class of models and not about one vendor.

\paragraph{On self-promotion as an attack surface.} The policy's text is inside
the judge's state, which makes ``write something the judge reads as evidence of
quality'' a gradient direction like any other. We both penalise it in the
rubric and track it independently, and we report the cross-reward matrix so
that a policy which moved towards its own judge is visible as a gap rather than
as an improvement.

\section{Limitations}
\jev{} is a closed model with no published architecture or weights, so its
calibration is an empirical property we measure rather than one we can reason
about from first principles, and it may shift between versions; we pin
\texttt{jev-1.13} and record the dated model string returned with every
response. Our policy scale (0.5B) is set by the compute available and the
ranking of reward sources may not transfer to substantially larger policies.
The prefix critic is probed at $K$ positions and interpolated, so the resulting
advantages are piecewise approximations of a true token-level value function.
Finally, using a hosted model's outputs as a training signal is subject to that
provider's terms of service, which we checked before running.

\section{Conclusion}
\todo{on completion of runs}

\bibliographystyle{plainnat}
\bibliography{refs}

\end{document}
